\begin{definition}[Unit filtration]\label{mod:localfield-unitfiltration}
Let $F$ be a nonarchimedean local field.  The subgroup
$\texttt{unitGroup}(F)\subset F^\times$ is Mathlib's canonical unit group of
the valuation subring, and its membership criterion is
\[
  u\in\texttt{unitGroup}(F)
  \quad\Longleftrightarrow\quad
  \operatorname{ord}_F(u)=0.
\]
It is canonically equivalent to the unit group of $\mathcal O_F$.  Define the
natural-number-indexed filtration with the manuscript's required case split
\[
 U_F^0=\mathcal O_F^\times,
 \qquad
 u\in U_F^{n+1}
 \quad\Longleftrightarrow\quad
 u\equiv 1\pmod{\mathfrak p_F^{n+1}}.
\]
Thus, only at positive depth,
\[
 u\in U_F^{n+1}
 \Longleftrightarrow
 \operatorname{CongruentAtDepth}_{n+1}(u,1)
 \Longleftrightarrow
 u-1\in\mathfrak p_F^{n+1}
 \Longleftrightarrow
 n+1\leq\operatorname{ord}_F(u-1).
\]
The first positive layer is exactly Mathlib's principal-unit group.

Every $U_F^n$ lies in $U_F^0$, and the filtration is antitone:
$m\leq n$ implies $U_F^n\subseteq U_F^m$.  It is closed under
multiplication and inversion; more generally, $u\in U_F^m$ and
$v\in U_F^n$ imply $uv\in U_F^{\min(m,n)}$.  For $u\in U_F^0$,
\[
 \operatorname{ord}_F(u^{-1}-1)=\operatorname{ord}_F(u-1),
\]
so inversion preserves every positive layer, while
$u\in U_F^{n+1}\setminus U_F^{n+2}$ is equivalent to
$\operatorname{ord}_F(u-1)=n+1$.  If
$x\in\mathfrak p_F^{n+1}$, the declaration \texttt{principalUnitOf}
constructs the unit $1+x\in U_F^{n+1}$.

Every unit-filtration layer is open in $F^\times$, and the positive layers
$\{U_F^{n+1}:n\in\mathbf N\}$ form a neighborhood basis of $1$ in
$F^\times$.  Thus any filtered family of units that eventually lies in
positive layers whose depths tend to infinity converges to $1$, both in
$F^\times$ and after coercion to $F$.

The multiplication error is computed exactly:
\[
 (uv-1)-\bigl((u-1)+(v-1)\bigr)=(u-1)(v-1).
\]
For $u\in U_F^{m+1}$ and $v\in U_F^{n+1}$ this error belongs to
$\mathfrak p_F^{m+n+2}$; in particular, for two elements of
$U_F^{n+1}$ it belongs to $\mathfrak p_F^{n+2}$.  This is the
multiplicative-to-additive linearization used on positive graded pieces.
For $u,v\in U_F^0$, division by $v$ gives
\[
 \operatorname{ord}_F(u/v-1)=\operatorname{ord}_F(u-v),
 \qquad
 u/v\in U_F^n
 \Longleftrightarrow
 \operatorname{CongruentAtDepth}_n(u,v).
\]

For certified indices $m\leq n$, \texttt{unitFiltrationInside} is the
copy of $U_F^n$ inside $U_F^m$, and its image back in $F^\times$ is exactly
$U_F^n$.  These internal denominators are nested at increasing depths.
The type
\[
 \texttt{UnitFiltrationQuotient}(F,m,n)=U_F^m/U_F^n
\]
is therefore available only with a proof of $m\leq n$.  Its canonical
representative map is surjective; two representatives have the same class
exactly when their ratio lies in $U_F^n$, equivalently when they are
congruent at depth $n$, and a representative has class one exactly when it
lies in $U_F^n$.  For $m\leq n_1\leq n_2$, there is a surjective canonical
projection $U_F^m/U_F^{n_2}\to U_F^m/U_F^{n_1}$ preserving representatives.

Finally,
\[
 \texttt{UnitGradedPiece}(F,n)=U_F^n/U_F^{n+1}.
\]
Its canonical representative map is surjective, with equality of classes
equivalent to depth-$(n+1)$ congruence and trivial class equivalent to
membership in $U_F^{n+1}$.  The degree-zero equivalence is the canonical
Mathlib equivalence
\[
 U_F^0/U_F^1\simeq (k_F)^\times.
\]
Writing
$\texttt{LatticeGradedPiece}(F,r)=\mathfrak p_F^r/\mathfrak p_F^{r+1}$,
the positive-depth map
\[
 U_F^{n+1}\longrightarrow
 \operatorname{Multiplicative}
   \bigl(\texttt{LatticeGradedPiece}(F,n+1)\bigr),
 \qquad u\longmapsto [u-1],
\]
is a surjective homomorphism with kernel $U_F^{n+2}$.  It induces the
equivalences
\[
 U_F^{n+1}/U_F^{n+2}
   \simeq_\times
   \operatorname{Multiplicative}
     (\mathfrak p_F^{n+1}/\mathfrak p_F^{n+2}),
 \qquad
 \operatorname{Additive}(U_F^{n+1}/U_F^{n+2})
   \simeq_+
   \mathfrak p_F^{n+1}/\mathfrak p_F^{n+2},
\]
and the multiplicative equivalence sends the class of $u$ to the class of
$u-1$.
\LeanFile{LanglandsFirstMainLemma/LocalField/UnitFiltration.lean}
\lean{LanglandsFirstMainLemma.unitGroup}
\lean{LanglandsFirstMainLemma.mem_unitGroup_iff_ord_eq_zero}
\lean{LanglandsFirstMainLemma.unitFiltration}
\lean{LanglandsFirstMainLemma.mem_unitFiltration_succ_iff_sub_mem_lattice}
\lean{LanglandsFirstMainLemma.unitFiltration_one_eq_principalUnitGroup}
\lean{LanglandsFirstMainLemma.unitFiltration_antitone}
\lean{LanglandsFirstMainLemma.unitFiltration_mul_mem_min}
\lean{LanglandsFirstMainLemma.ord_inv_sub_one_eq}
\lean{LanglandsFirstMainLemma.mem_unitFiltration_and_not_mem_succ_iff}
\lean{LanglandsFirstMainLemma.principalUnitOf}
\lean{LanglandsFirstMainLemma.unitFiltration_isOpen}
\lean{LanglandsFirstMainLemma.unitFiltration_hasBasis_nhds_one}
\lean{LanglandsFirstMainLemma.tendsto_one_of_eventually_mem_unitFiltration}
\lean{LanglandsFirstMainLemma.tendsto_coe_one_of_eventually_mem_unitFiltration}
\lean{LanglandsFirstMainLemma.unitFiltration_mul_error_mem}
\lean{LanglandsFirstMainLemma.unitFiltration_mul_error_mem_next}
\lean{LanglandsFirstMainLemma.div_mem_unitFiltration_iff_congruentAtDepth}
\lean{LanglandsFirstMainLemma.unitFiltrationInside}
\lean{LanglandsFirstMainLemma.UnitFiltrationQuotient}
\lean{LanglandsFirstMainLemma.unitFiltrationQuotientMk_eq_mk_iff_congruentAtDepth}
\lean{LanglandsFirstMainLemma.unitFiltrationQuotientProjection}
\lean{LanglandsFirstMainLemma.UnitGradedPiece}
\lean{LanglandsFirstMainLemma.unitGradedZeroEquivResidueFieldUnits}
\lean{LanglandsFirstMainLemma.positiveUnitToLatticeGraded}
\lean{LanglandsFirstMainLemma.positiveUnitToLatticeGraded_ker}
\lean{LanglandsFirstMainLemma.positiveUnitGradedEquivLattice}
\lean{LanglandsFirstMainLemma.positiveUnitGradedAddEquivLattice}
\leanok
\SourceText{Local notation; Theorem thm:norm-filtration.}
\uses{mod:localfield-lattices}
\FormalizationContract
\end{definition}
